\documentclass[aspectratio=169,10pt]{beamer}

\usetheme{NMA}
\usepackage{booktabs}
\usepackage{tabularx}
\usepackage{array}
\newcolumntype{Y}{>{\raggedright\arraybackslash}X}

\title[NMA Project Overview]{Project Overview}
%\subtitle{Diagnosing where neural processing diverges on miss trials}
\author[C. Lin]{Presented by Chunyi Lin}

\date{}
\course{Banshee Fireweed Group 1}
\week{Week 2}
\session{Day 4}

\newcommand{\sourceNote}[1]{%
  \par\vfill{\fontsize{6}{7}\selectfont\color{NMAMidGray}#1}%
}

\begin{document}

\begin{frame}[plain]
  \titlepage
\end{frame}

\begin{frame}{The same change of stimuli produces two actions}

  \begin{columns}[c,onlytextwidth]
    \begin{column}{.46\textwidth}
      \begin{block}{Hit}
        \centering\Large visual change $\rightarrow$ lick
      \end{block}
    \end{column}
    \begin{column}{.46\textwidth}
      \begin{alertblock}{Miss}
        \centering\Large visual change $\rightarrow$ no lick
      \end{alertblock}
    \end{column}
  \end{columns}

  \vspace{1em}
  The external event is similar, but the behavioral output differs. This makes
  miss trials a natural probe of variability generated inside the animal.

  %\NMAtakeaway{The question is not whether hits and misses differ, but where their processing pathways diverge.}
\end{frame}

\begin{frame}{A miss can arise at four different stages}
  \centering
  {\Large
    internal state $\longrightarrow$ sensory encoding
    $\longrightarrow$ decision $\longrightarrow$ action
  }

  \vspace{1.1em}
  \begin{columns}[T,onlytextwidth]
    \begin{column}{.48\textwidth}
      \begin{itemize}
        \item The trial begins in a different population state.
        \item The visual change is encoded less effectively.
      \end{itemize}
    \end{column}
    \begin{column}{.48\textwidth}
      \begin{itemize}
        \item Sensory evidence fails to drive a decision.
        \item A decision forms but fails to produce licking.
      \end{itemize}
    \end{column}
  \end{columns}

  \NMAtakeaway{A single hit--miss decoder cannot identify which transformation failed.}
\end{frame}

\begin{frame}{The experiment resolves each trial}
  \begin{columns}[T,onlytextwidth]
    \begin{column}{.48\textwidth}
      \textbf{Comparable task event}

      Hits and misses both come from change trials, reducing the largest
      stimulus difference.

      \vspace{.8em}
      \textbf{Population activity}

      Simultaneous neurons provide a high-dimensional view of cortical state.
    \end{column}
    \begin{column}{.48\textwidth}
      \textbf{Resolved behavioral outcome}

      Licking identifies the response; reaction time bounds activity before
      overt action.

      \vspace{.8em}
      \textbf{Arousal and movement}

      Pupil, running, and licking provide competing explanations for neural
      differences.
    \end{column}
  \end{columns}

  \sourceNote{Pupil controls are motivated by Reimer et al., 2016, Nature Communications.}
\end{frame}

\begin{frame}{A miss is an outcome, not an internal state}
  \begin{columns}[c,onlytextwidth]
    \begin{column}{.46\textwidth}
      \begin{alertblock}{Avoid this shortcut}
        \[
          \text{miss}=\text{disengaged}
        \]
        An engaged mouse can miss, and a weakly engaged mouse can occasionally
        respond correctly.
      \end{alertblock}
    \end{column}
    \begin{column}{.48\textwidth}
      \begin{block}{Infer state from the sequence}
        \[
          P(z_i=\text{engaged}\mid\text{behavioral history})
        \]
        Persistent state is a hypothesis about temporal structure, not a label
        assigned from one trial.
      \end{block}
    \end{column}
  \end{columns}

  \NMAtakeaway{Model behavior first; then ask how inferred state appears in neural activity.}
\end{frame}

\begin{frame}{Do failures cluster into persistent modes?}
  \begin{tabularx}{\textwidth}{>{\bfseries\raggedright\arraybackslash}p{.22\textwidth}Y Y}
    \toprule
    Model & What it captures & Key test \\
    \midrule
    Fixed lapse & Independent stimulus-unrelated errors & Is one lapse rate sufficient? \\
    History GLM & Continuous dependence on recent trials & Does history explain clustering? \\
    GLM--HMM & Distinct, persistent stimulus--response mappings & Do latent states improve held-out prediction? \\
    \bottomrule
  \end{tabularx}

  \vspace{.8em}
  The output is a posterior state probability for every trial, not a binary
  relabeling of misses.

  \sourceNote{Ashwood et al., 2022, Nature Neuroscience.}
\end{frame}

\begin{frame}{Match sensory evidence before decoding}
  \begin{columns}[T,onlytextwidth]
    \begin{column}{.52\textwidth}
      Match hit and miss trials within:
      \begin{itemize}
        \item previous and current image identity;
        \item image-transition type;
        \item recent stimulus and adaptation history;
        \item trial position or temporal block;
        \item relevant task timing.
      \end{itemize}
    \end{column}
    \begin{column}{.42\textwidth}
      \begin{alertblock}{Without matching}
        A decoder may separate
        \[
          \text{easy} \quad \text{vs.} \quad \text{difficult changes}
        \]
        rather than internal states.
      \end{alertblock}
    \end{column}
  \end{columns}

  \NMAtakeaway{``Both are change trials'' does not guarantee equivalent sensory input.}
\end{frame}

\NMAsectionframe[Define each processing stage independently]{Four neural axes localize the failure}

\begin{frame}{A state signal should exist before the change}
  \begin{columns}[c,onlytextwidth]
    \begin{column}{.48\textwidth}
      Use pre-change population activity to predict:
      \begin{itemize}
        \item behaviorally inferred latent state;
        \item current-trial performance;
        \item next-trial performance.
      \end{itemize}
    \end{column}
    \begin{column}{.46\textwidth}
      \begin{block}{Pre-change state axis}
        \[
          \mathbf{x}_{\mathrm{pre}}
          \longrightarrow P(\mathrm{engaged})
        \]
      \end{block}
      Do not subtract each trial's baseline: tonic pre-change differences are
      the signal of interest.
    \end{column}
  \end{columns}

  \NMAtakeaway{Predicting future behavior rules out current licking or reward as the cause.}
\end{frame}

\begin{frame}{Define sensory coding without outcome}
  \begin{columns}[T,onlytextwidth]
    \begin{column}{.48\textwidth}
      Train a sensory decoder on a variable such as:
      \begin{itemize}
        \item change versus no change;
        \item image identity;
        \item changed versus repeated image;
        \item stimulus transition.
      \end{itemize}
      Then project matched hits and misses onto this axis.
    \end{column}
    \begin{column}{.47\textwidth}
      \begin{block}{Two distinct outcomes}
        \textbf{Weak signal on misses}

        Sensory encoding may have failed.

        \vspace{.65em}
        \textbf{Normal signal on misses}

        The failure lies downstream.
      \end{block}
    \end{column}
  \end{columns}
\end{frame}

\begin{frame}{Choice information must precede the lick}
  Train a cross-validated hit--miss classifier on early post-change,
  pre-lick population activity.

  \vspace{.8em}
  \begin{columns}[T,onlytextwidth]
    \begin{column}{.48\textwidth}
      \begin{block}{What it establishes}
        Early population activity contains reproducible information about the
        eventual behavioral response.
      \end{block}
    \end{column}
    \begin{column}{.48\textwidth}
      \begin{alertblock}{What it does not establish}
        A supervised score is not automatically an accumulator, drift rate, or
        mechanistic decision variable.
      \end{alertblock}
    \end{column}
  \end{columns}

  \NMAtakeaway{Call it a choice-related population score until accumulator-like properties are tested.}
\end{frame}

\begin{frame}{A motor axis isolates impending action}
  Define lick-related activity using trials or events not used for the
  hit--miss axis:
  \begin{itemize}
    \item spontaneous licks or separate lick bouts;
    \item false alarms;
    \item held-out pre-lick activity;
    \item other movement measurements when available.
  \end{itemize}

  \begin{columns}[T,onlytextwidth]
    \begin{column}{.47\textwidth}
      \begin{block}{Sensory present, choice absent}
        Evidence was encoded but did not reach a decision-related response.
      \end{block}
    \end{column}
    \begin{column}{.47\textwidth}
      \begin{exampleblock}{Choice present, motor absent}
        A decision-related response emerged but did not produce licking.
      \end{exampleblock}
    \end{column}
  \end{columns}
\end{frame}

\begin{frame}{A gate changes sensory-to-choice coupling}
  \[
    C_i=\beta_0+\beta_1S_i+\beta_2Z_i
      +\beta_3(S_i\times Z_i)+\varepsilon_i
  \]

  \begin{itemize}
    \item $S_i$: sensory-axis activity
    \item $C_i$: choice-axis activity
    \item $Z_i$: engagement-state probability
  \end{itemize}

  \begin{block}{Mechanistic test}
    The interaction $\beta_3$ asks whether the mapping from sensory evidence to
    choice-related activity changes across internal states.
  \end{block}

  \NMAtakeaway{Gating is a changed transformation, not merely a weaker response.}
\end{frame}

\NMAsectionframe[Separate internal state from competing explanations]{Validation gives the axes meaning}

\begin{frame}{Controls must compete with neural state}
  \begin{align*}
    M_0:\quad \text{behavior} &\sim
      \text{pupil}+\text{running}+\text{movement}+\text{history} \\
    M_1:\quad \text{behavior} &\sim
      \text{pupil}+\text{running}+\text{movement}+\text{history}
      +\text{neural state}
  \end{align*}

  Fit every nuisance transformation on training trials only, then compare
  held-out performance.

  \begin{alertblock}{Calibrated conclusion}
    Neural population activity predicts engagement beyond the measured arousal
    and movement variables. It does not prove independence from arousal.
  \end{alertblock}

  \sourceNote{Reimer et al., 2016; Ortiz et al., 2020.}
\end{frame}

\begin{frame}{Slow time can masquerade as internal state}
  Motivation, neural activity, performance, and recording quality can all drift
  across a session.

  \vspace{.6em}
  Use:
  \begin{itemize}
    \item time-matched hit--miss comparisons;
    \item blocked or leave-one-block-out cross-validation;
    \item within-session permutations that preserve local temporal structure;
    \item explicit sensitivity analyses for slow drift.
  \end{itemize}

  \NMAtakeaway{The model must generalize beyond recognizing early versus late session.}
\end{frame}

\begin{frame}{The axes turn every miss into a diagnosis}
  \renewcommand{\arraystretch}{1.25}
  \begin{tabularx}{\textwidth}{>{\bfseries\raggedright\arraybackslash}p{.32\textwidth}Y}
    \toprule
    Held-out result & Most direct interpretation \\
    \midrule
    Sensory signal is weak & The visual change was encoded less effectively. \\
    Sensory is preserved; choice is weak & Sensory evidence failed to drive a decision-related response. \\
    Sensory and choice are preserved; motor is weak & A decision-related response failed to produce action. \\
    Pre-change state predicts current and future behavior & Outcomes depend on a persistent internal condition. \\
    \bottomrule
  \end{tabularx}

  \vspace{.6em}
  These interpretations require matched inputs, independent axes, and blocked
  validation.
\end{frame}

\begin{frame}{Estimate within sessions before pooling}
  Neural axes live in different neuron spaces across recordings, and trials are
  nested within sessions and animals.

  \vspace{.8em}
  \begin{columns}[T,onlytextwidth]
    \begin{column}{.47\textwidth}
      \begin{alertblock}{Avoid}
        Pooling all trials as independent observations, which inflates precision
        and hides heterogeneous sessions.
      \end{alertblock}
    \end{column}
    \begin{column}{.47\textwidth}
      \begin{block}{Prefer}
        Estimate held-out effects per session, then combine them with a
        hierarchical or mixed-effects model.
      \end{block}
    \end{column}
  \end{columns}

  \NMAtakeaway{The population claim is recurrence across sessions or animals.}
\end{frame}

\begin{frame}{Prior work motivates the test, not the answer}
  \begin{itemize}
    \item \textbf{Lapses can be persistent.} Mouse choices can occupy sensory-guided
      and biased latent strategies for tens to hundreds of trials.
    \item \textbf{Engagement is not simply thirst or pupil size.} Performance and
      V1 responses can decline without a matching decrease in pupil diameter.
    \item \textbf{Strategy, engagement, and cortical state differ.} A neural
      hit--miss axis may reflect strategy, expectation, movement, engagement, or
      their interaction.
  \end{itemize}

  \NMAtakeaway{The analysis must distinguish these explanations rather than assign an engagement label by default.}

  \sourceNote{Ashwood et al., 2022; Ortiz et al., 2020; Jacobs et al., 2020; Piet et al., 2024.}
\end{frame}

\begin{frame}{Move from behavior to mechanism}
  \begin{enumerate}
    \item Infer persistent behavioral state.
    \item Match the sensory and temporal conditions.
    \item Define pre-change, sensory, choice, and motor axes independently.
    \item Test state-dependent sensory-to-choice coupling.
    \item Validate against movement, arousal, time, and session structure.
  \end{enumerate}

  \vspace{.7em}
  \begin{block}{Role of the current notebook}
    The existing hit--miss decoder is the \textbf{choice-related axis module}.
    It is one component of the internal-state analysis, not the complete model.
  \end{block}
\end{frame}

\begin{frame}{References}
  \begin{columns}[T,onlytextwidth]
    \begin{column}{.48\textwidth}
      {\scriptsize
      \begin{itemize}
        \item Wichmann \& Hill (2001). The psychometric function I.
        \item Ashwood et al. (2022). Mice alternate between discrete strategies during perceptual decision-making.
        \item Ortiz et al. (2020). Motivation and engagement during visually guided behavior.
        \item Reimer et al. (2016). Pupil fluctuations track rapid changes in adrenergic and cholinergic activity.
      \end{itemize}}
    \end{column}
    \begin{column}{.48\textwidth}
      {\scriptsize
      \begin{itemize}
        \item Jacobs et al. (2020). Cortical state fluctuations during sensory decision making.
        \item Piet et al. (2024). Behavioral strategy shapes activation of the Vip--Sst inhibitory circuit.
        \item Full links and proposal context are preserved in the source narrative supplied with this deck.
      \end{itemize}}
    \end{column}
  \end{columns}
\end{frame}

\NMAclosingframe[Was the change not encoded, or was it prevented from reaching action?]{A miss reveals the failed stage}

\end{document}
